\SourceLeaf{33} Mitteilung von Herrn \textsc{Mehmke} verdanke ich die Kenntnis, dass \textsc{Gauss} auch in der
zweiten Tabelle (abgesehen von den drei genannten Fällen) für das System
derjenigen Einheiten, die sich mithilfe ganzzahliger $a$, $b$, $c$ in der
Gestalt $a+b\nu+c\nu\nu$ darstellen lassen, jedes Mal die zur
Fundamentaleinheit reziproke Einheit angegeben hat. In der Mehrzahl der
Fälle handelt es sich auch um die Fundamentaleinheit des Körpers.
Ausnahmen von dieser Regel liegen vor bei $n=16$, wo die achte Potenz der
Fundamentaleinheit angegeben ist, und bei $n=10$ und $n=19$, wo die
reziproken Fundamentaleinheiten die folgenden sind.

\[
\frac{7+\sqrt[3]{10}-2\sqrt[3]{100}}{3}, \qquad
\frac{2+2\sqrt[3]{19}-\sqrt[3]{19^2}}{3}.
\]

\par\noindent\hfill\textsc{Fricke.}\par\endgroup
